\section{The decks of Grothendieck's notes}\label{sec:grdecks}

Three decks recur in Grothendieck's working notes. Here is what the three arches of \S\ref{sec:morearches} say about them (Lean: \lean{GrothendieckDecks.lean}).
\begin{itemize}
\item The simplex category $\Delta$, the first test category of \emph{Pursuing Stacks} (folders 108 and 118 of the fonds): simplicial sets form a two-valued and connected topos, since there is a constant map between any two objects of $\Delta$; this topos is not the presheaf topos of a preorder, since the two vertices of $[1]$ are two maps $[0]\to[1]$ (\lean{simplexCategory_twoValued}, \lean{simplexCategory_isConnected}, \lean{simplexCategory_not_spatial}).
\item A group $G$ as a one-object groupoid, whose presheaf topos is $\B G$ (folder 161-6, p.~22, where Grothendieck reduces a question to this case): the topos is two-valued (as \cite[IV~8.4.3]{sga4} notes for any monoid) and connected, and it is the presheaf topos of a preorder iff $G$ is trivial (\lean{singleObj_twoValued}, \lean{singleObj_isConnected}, \lean{singleObj_spatial_iff}).
\item A preordered set $I$, whose presheaf topos $\widehat I$ is $\mathrm{Top}(I)$ (folder 37, p.~48; folder 161-3, p.~26): the topos is two-valued iff any two elements are comparable both ways, that is, iff its posetal quotient has at most one point (\lean{preorder_twoValued_iff}).
\end{itemize}
These are immediate corollaries. We record them because they are the first entries of the registry taken from the fonds.
